This section defines the common execution contract used by the instruction descriptions. An individual
instruction may add stricter rules, but it does not silently replace the rules below. The purpose of this
section is to make operand evaluation, side effects, repetition, and status updates readable before the
per-instruction pages.

\subsection{Instruction Boundary and Default Execution Rules}
The instruction boundary is fixed by word 0 before operand evaluation begins. An implementation may reject an
instruction before any architectural state is changed if the encoded length cannot contain the selected form.

@DEFAULT_RULE_TABLE@

\subsection{Instruction Boundary, Operand Evaluation, and EA Defaults}
\begin{itemize}
\item Word 0 length selects the instruction boundary; prefixes and extension words never extend it implicitly.
\item Operands are decoded in instruction order; source reads complete before the final destination write unless an atomic form says otherwise.
\item EA calculation may produce an address without reading memory; memory is read only when the operation needs the operand value.
\end{itemize}

\subsection{Mnemonic Suffix Rules}
Conditional mnemonic forms use the condition names listed in the Condition Codes section.

@SUFFIX_RULE_TABLE@

\subsection{Operand and Status Notation}
The generated operation fields use compact notation. The following terms define when an operand is only
addressed, when it is read, and which status register is affected.

@NOTATION_TABLE@

\subsection{Shared Side-Effect Rules}
These rows summarize execution rules that apply to whole instruction families. The instruction descriptions
list exact forms and operands; this block records the common side-effect model.

@SHARED_SIDE_EFFECT_BLOCK@
